\documentclass{article}
\usepackage{propositions}

%% align=nextline and align=flush-nextline give the label a line of its own and
%% start the body on the next one.  Starting it is the delicate part: \noindent
%% would begin a paragraph there and then, which suits a body of text but not
%% one that opens with vertical material.  A nested prop's \begin{list} closes
%% such a paragraph at once, and an empty paragraph is a line of white space --
%% so the sublist came out separated from its label by a blank line, while the
%% same label followed by text was not.  The body is therefore started with
%% \@afterheading, which unindents whichever paragraph comes next rather than
%% opening one now.

\newcommand\gap{\rule[-2pt]{0.4pt}{8pt}}

\begin{document}

\section{A nextline label followed by a list}

The first two lists below must look the same below their labels: the body
begins on the very next line either way, whether it is a sublist or a
paragraph.  The third shows what \verb|nextline| is for --- without it the
sublist starts on the label's own line and is pushed right, which is what makes
its internal alignment go awry.

\noindent\rule{\textwidth}{0.3pt}
\begin{prop}
  \item[Label A, align=nextline]
  \begin{prop}
    \item First sub-item.
    \item Second sub-item.
  \end{prop}
\end{prop}
\noindent\rule{\textwidth}{0.3pt}
\begin{prop}
  \item[Label B, align=nextline] Ordinary text after the label, which must
  begin on the next line exactly as the sublist above does.
\end{prop}
\noindent\rule{\textwidth}{0.3pt}
\begin{prop}
  \item[Label C]
  \begin{prop}
    \item First sub-item, on the label's own line (no \texttt{nextline}).
    \item Second sub-item.
  \end{prop}
\end{prop}
\noindent\rule{\textwidth}{0.3pt}

Measured on the baselines: A and B must both advance one \verb|\baselineskip|
from label to body (11.95pt in \texttt{article} at 10pt), and C must advance
none, its sublist sharing the label's line.

\section{The same for \texttt{flush-nextline}}

\noindent\rule{\textwidth}{0.3pt}
\begin{prop}
  \item[Label D, align=flush-nextline]
  \begin{prop}
    \item Sub-item under a flush-nextline label.
  \end{prop}
\end{prop}
\noindent\rule{\textwidth}{0.3pt}

\section{Nested nextline labels}

Each level must advance exactly one line, with no blank line accumulating as
the nesting deepens.

\noindent\rule{\textwidth}{0.3pt}
\begin{prop}
  \item[Outer, align=nextline]
  \begin{prop}
    \item[Inner, align=nextline]
    \begin{prop}
      \item Innermost.
    \end{prop}
  \end{prop}
\end{prop}
\noindent\rule{\textwidth}{0.3pt}

\section{An empty body, and a following paragraph}

A \verb|nextline| label with nothing after it must not disturb the item that
follows, nor leave anything behind that unindents a later paragraph: the
\verb|\everypar| hook \verb|\@afterheading| leaves pending is cleared by the
paragraphs it fires on, and is confined to the list's group in any case.

\noindent\rule{\textwidth}{0.3pt}
\begin{prop}
  \item[Empty body, align=nextline]
  \item[Next item] This item must be laid out normally.
\end{prop}
\noindent\rule{\textwidth}{0.3pt}

This paragraph follows the list and must be indented like any other opening
paragraph of running text, long enough to wrap so that the first line's indent
can be compared with the second's.

\end{document}
